% TOP_PROBLEM: {"id": 1124, "title": "Kakeya Needle Problem", "category_id": 6, "category": {"id": 6, "name": "geometry", "display_name": "Geometry", "description": "Euclidean and non-Euclidean geometry, geometric structures.", "slug": "geometry", "order_index": 6, "created_at": "2026-07-31T15:26:25.671Z"}, "status": "open"}
% FIRST_POSTED: null
% ENTRY_KIND: statement_only
% Imported from: batch04.tex; UnsolvedMath ID 1124
\documentclass[11pt]{article}
\usepackage[margin=25mm]{geometry}
\usepackage{fontspec}
\setmainfont{Latin Modern Roman}
\usepackage{amsmath,amssymb,mathtools}
\usepackage{unicode-math}
\setmathfont{Latin Modern Math}
\usepackage{xurl,hyperref}
\hypersetup{colorlinks=true,urlcolor=blue}
\setcounter{secnumdepth}{0}
\setlength{\parindent}{0pt}
\setlength{\parskip}{5pt}
\setlength{\emergencystretch}{3em}
\newcommand{\sref}[2]{\hyperlink{source:#1:#2}{[S#2]}}
\newcommand{\cataloguescope}[1]{\paragraph{Recorded target.}#1}
\begin{document}
% UnsolvedMath ID 1124; source catalogue code GEO-004
\setcounter{section}{1123}
\section{Kakeya Needle Problem}\label{1124}
{\small\sffamily UnsolvedMath ID 1124 · Geometry}
\subsection{Definitions and mathematical statement}

\paragraph{Shared notation.}$\mathbb N=\{1,2,\ldots\}$, and $\mathbb Z,\mathbb Q,\mathbb R,\mathbb C$ denote the integers, rationals, reals and complex numbers. Any other conventions are specified locally.


Let \(K\subseteq\mathbb R^2\) be a bounded Lebesgue-measurable region.
A closed full rotation of a unit needle within \(K\) is described by
continuous functions \(a:[0,1]\to\mathbb R^2\) and
\(\theta:[0,1]\to\mathbb R\) satisfying
\[
 a(1)=a(0),\qquad \theta(1)=\theta(0)+2\pi,\qquad
 \{a(t)+u(\cos\theta(t),\sin\theta(t)):0\le u\le1\}
       \subseteq K\quad (0\le t\le1).
\]
Translations during rotation are allowed; neither convexity nor simple
connectedness is imposed.  Let \(\mathcal K\) be the class of such regions
and write \(|K|\) for planar Lebesgue area.
\paragraph{Statement recorded in the supplied source.}
Determine
\[
 \inf_{K\in\mathcal K}|K|
\]
and whether this infimum is attained.  This asks for the least area
needed for an actual continuous \(360^\circ\) rotation returning the
needle to its starting position, rather than merely a set containing
a unit segment in each direction. \sref{1124}{2} \sref{1124}{3}
\subsection{Short English statement}
How little planar area can permit a unit needle to turn continuously through a full rotation and return to its starting position? Is there a region of minimum area?
\subsection{Sources}
\begin{description}
\item[\hypertarget{source:1124:1}{S1}] ULAM.AI and the UnsolvedMath Contributors, \emph{UnsolvedMath: A Curated Collection of Open Mathematics Problems}, record 1124 (GEO-004). CC BY 4.0. \url{https://huggingface.co/datasets/ulamai/UnsolvedMath}.
\item[\hypertarget{source:1124:2}{S2}] Hee-Kap Ahn, Sang Won Bae, Otfried Cheong, Joachim Gudmundsson, Takeshi Tokuyama and Antoine Vigneron, \emph{A Generalization of the Convex Kakeya Problem}, arXiv:1209.2171, Section 1, discussion of the unrestricted problem on pages 1--2. \url{https://arxiv.org/abs/1209.2171}.
\item[\hypertarget{source:1124:3}{S3}] Terence Tao, \emph{Rotating needles in space: the road to the Kakeya conjecture, and why it matters}, arXiv:2608.22209 (2026), Section 1, Problem 1.1 and footnotes 1--2. \url{https://arxiv.org/abs/2608.22209}.
\end{description}
\label{end:1124}
\end{document}
